\subsubsection{Justificación de la solución [Pendientes: 1. sustentar la comparación con las alternativas de la sección 5.2]}\label{sec:justificacion}
La solución se justifica por la posibilidad de separar el esfuerzo de preparación científica del momento en que se necesita consultar información para decidir. En lugar de exigir al usuario configurar modelos y esperar sus ejecuciones, el producto busca convertir resultados previamente organizados en una representación interpretable de la evolución posible del incendio. El valor esperado se encuentra en esa articulación de biblioteca, consulta y comunicación espacial, mientras que sus fundamentos computacionales se desarrollan en la sección~\ref{sec:fundamentos}. Este enfoque resulta pertinente para la necesidad expresada en las entrevistas, aunque su utilidad depende de que las condiciones recuperadas representen suficientemente la situación consultada y de que la interfaz permita entender sus limitaciones. En consecuencia, la propuesta incluye comprobaciones técnicas y evaluación de comprensión, en vez de atribuir por anticipado exactitud física o efectividad operativa a una animación.

La elección también busca acotar un problema territorial y computacional amplio a una contribución realizable de software, con preparación trazable de datos y resultados. Sin embargo, este argumento de pertinencia aún no constituye una comparación completa de soluciones, dado que las tres alternativas adicionales exigidas por la plantilla no se han establecido. La justificación comparativa se completará a partir del análisis de la sección~\ref{sec:alternativas}, utilizando criterios comunes de alcance, datos, costo de preparación, consulta y presentación de resultados. Hasta entonces, no se afirma que la propuesta supere a herramientas existentes ni que las tesis empleadas como antecedentes equivalgan, por sí mismas, a alternativas integrales para el mismo problema.
